% Discusión, limitaciones, aplicación a otras redes y problemas abiertos (memoria, capítulos 9 y 10).

\subsection{Where the accuracy comes from}

Three observations separate the contributions. First, knowing which segment the vehicle is in carries most of the
information: deleting every GPS position, so that only the stop registered by the feed remains, costs nothing below
three minutes and three points at ten to twenty minutes (Table~\ref{tab:ablation}). What matters, as
Proposition~\ref{prop:differences} anticipates, is to start from where the vehicle is rather than from the published
time. Second, correcting the shape of the timetable already improves on it a lot: without any live data, the corrected
timetable doubles the hit rate of the published one. Third, the coverage of learned segments is what moves accuracy the
most: in the technical report's sample, arrivals under ten minutes were within $\pm2$~minutes in 93\% of cases when the
whole path was learned, 76\% with one unlearned segment and 59\% with two or more. The remaining error is limited by the
data themselves, one position per minute and a ground truth known to $\pm30$~s, and by what no model can foresee with
them: traffic lights, traffic and dwell times.

A neural network that also uses what happens that day reduces the error of segment sums by 6--7\%, and a Kalman filter
on the vehicle's own pace by 4\% with two parameters. The method keeps medians per segment, a shrinkage with a single
constant and a four-level fallback because every parameter has a physical meaning and can be checked by hand, because
any figure the app shows can be explained, and because it runs without a server. The own-pace filter is the natural
next step, once it is measured on the full arrival.

\subsection{Abstention}

Several rules do not increase the measured hit rate but avoid the errors that most damage trust: labelling estimates
as estimates, hiding a delay that is not credible while keeping the arrival, distrusting two sources that contradict
each other, and saying that a bus may already have passed rather than announcing it as coming. On 3 November 2026,
97.8\% of the trips shown as \texttt{likely\_departed} were not seen in the feed for the rest of the day.

\subsection{Limitations}
\label{sec:limitations}

\begin{itemize}
  \item \textbf{A changing method.} The figures describe the app as it ran each day, with calibrations and rules that
        changed during the evaluation, not a frozen version.
  \item \textbf{Ground truth.} The actual passage is known to $\pm30$~s, which especially affects the $\pm1$~minute
        figures, and it is the moment the feed changes its next stop, which may lag slightly behind the physical passage
        (equally for the method and the timetable).
  \item \textbf{Censoring and dependence.} Only passages observed within the capture are evaluated, so very late buses
        may be missed; successive predictions for the same vehicle are correlated, so the effective sample is smaller
        than the number of predictions.
  \item \textbf{Season and holidays.} The data cover October and early November 2026, with a single weekday holiday;
        summer and school holidays are unmeasured.
  \item \textbf{What the feed does not say.} A trip without a vehicle in the feed is shown with layer 3 even if it does
        not run (about 4\% of the trips that should be running cannot be interpreted at the moment), extra services
        not in the timetable do not appear, and night trips are not measured because the feed is frozen.
  \item \textbf{One network, one evaluator.} All results are from Bizkaibus and were obtained by the author with his own
        tools. The reference implementation, the test vectors and the public evaluator allow checking that the method
        described is the one measured, but they do not replace an independent replication on other data.
\end{itemize}

\subsection{Applying the method to another network}
\label{sec:other-networks}

The method needs, in order of importance: a GTFS timetable; a real-time feed that identifies the trip of each vehicle
and its next stop, with the time of the record (enough for layers 2 and 3); the GPS position and its time (for layer 1
and learning); shapes with distances (otherwise localisation uses straight lines between stops); and a way to capture
the feed periodically. Only the reading of the feed and the matching depend on the operator's format. The time bands
reflect demand in Biscay, and four constants were measured there and should be re-measured: the 50~m corridor, the
20-minute limit of layer 2, the 45-minute cap and the 60~s late departure from the terminal. With no samples, the model
reduces to timetable differences with $k=1$, which already removes the accumulated delay; learned segments appear after
a few days of capture. The steps for another city are: check what the operator publishes; build the timetable database
with merged variants, aliases and corrected shapes; adapt the reading of the feed and check that every vehicle finds
its trip; start the capture and the daily recalibration; measure with the protocol of Section~\ref{sec:protocol} before
showing anything as live; and check that the implementation reproduces the test vectors. If the operator already
publishes predictions per stop (GTFS-Realtime Trip Updates), they can be used directly after checking that they are
real data and not a copy of the timetable.

\subsection{Open problems}
\label{sec:open}

The following questions are open and suitable for further research, including bachelor's and master's theses; the
public data and evaluator allow comparing any answer with the method on equal terms.

\begin{enumerate}
  \item \textbf{Own pace.} Measure the Kalman filter on the vehicle's own pace on the full arrival shown to the user.
  \item \textbf{Prediction intervals.} Estimate the dispersion of each segment, not only its median, and show an
        interval (``in 6 to 8 minutes'') with a measured coverage.
  \item \textbf{Comparison with the operator's predictions.} Compare the method with the waiting times that the operator
        itself shows, on the same vehicles and stops; a capture for this purpose is running in the public repository.
  \item \textbf{Season, special days and events.} With months of data, measure whether summer, school holidays,
        festivities or events should be learned separately.
  \item \textbf{Trips without a vehicle.} Distinguish a trip that should already have left and has no vehicle in the
        feed from one that is merely late.
  \item \textbf{Vehicle sequences.} Follow the vehicle on its previous trip to estimate the real departure from the
        terminal, if the assignment of vehicles to trips repeats from day to day.
  \item \textbf{Other networks.} Apply the method to other cities with position-only feeds and check how much of what
        was measured here holds.
  \item \textbf{Journey planning.} Feed the predicted arrivals to a planner with transfers so that proposed journeys use
        real times rather than published ones.
\end{enumerate}
